\documentclass[tikz]{standalone}
\usepackage{amsmath}
\definecolor{bgcolor}{HTML}{171D2B}
\definecolor{aksen}{HTML}{F7B538}
\begin{document}
\pagecolor{bgcolor}
\begin{tikzpicture}[xscale=1.5, yscale=1.0]
  \def\N{5}
  % koefisien binomial C(n,k), dihitung manual
  \foreach \n/\vals in {0/{1}, 1/{1,1}, 2/{1,2,1}, 3/{1,3,3,1}, 4/{1,4,6,4,1}, 5/{1,5,10,10,5,1}} {
    \foreach \k [count=\i from 0] in \vals {
      \node[text=white, font=\large] (n\n-\i) at ({(\i-\n/2)*1.5}, {-\n*1.0}) {\k};
    }
  }
  % garis penopang: setiap angka dihubungkan ke dua angka di atasnya
  \foreach \n in {1,...,5} {
    \foreach \k in {0,...,\n} {
      \pgfmathtruncatemacro{\l}{\k-1}
      \ifnum\k>0 \draw[gray!60, thin] (n\the\numexpr\n-1\relax-\l) -- (n\n-\k); \fi
      \pgfmathtruncatemacro{\r}{\k}
      \ifnum\k<\n \draw[gray!60, thin] (n\the\numexpr\n-1\relax-\r) -- (n\n-\k); \fi
    }
  }
  % sorot satu contoh penjumlahan: 6 = 3 + 3 (baris 4, kolom 2)
  \node[text=aksen, font=\large] at (n4-2) {};
  \draw[aksen, thick, rounded corners] ([xshift=-3mm,yshift=3mm]n3-1.north west) rectangle ([xshift=3mm,yshift=-3mm]n3-2.south east);
  \draw[aksen, thick, rounded corners] ([xshift=-3mm,yshift=3mm]n4-2.north west) rectangle ([xshift=3mm,yshift=-3mm]n4-2.south east);
\end{tikzpicture}
\end{document}
